\par\clearpage
\noindent\makebox[0pt][l]{\fontsize{13}{16}\selectfont \textbf{382}}\hfill{\small\textbf{CORRESPONDENCE WITH GERLING.}}\hfill\null\par\medskip
\begin{center}
\textsc{Gauss} to \textsc{Gerling}.\quad Göttingen, 27~July 1823.
\end{center}

\par\noindent\hspace*{\parindent}\mbox{.\ .\ .\ .\ .\ .}\quad I have had \textsc{Repsold} make me another new eyepiece for my theodolite, with spider threads of almost incredible fineness, \(29^{\prime\prime}\) apart.
(Those in \textsc{Ertel}'s eyepiece, which are already quite fine but much thicker than those just mentioned, are \(36^{\prime\prime}\) apart.)
This arrangement is extremely advantageous in observing distant \textls[120]{faint} terrestrial objects. Unfortunately I received it only post festum.
The Hamburg tower, which almost always looked very faint from Nindorf, Timpenberg and Lüneburg because of the smoky haze in the air, gave me a great deal of trouble, as did the Lüneburg tower from Hamburg. On the whole the measurements in those parts are not quite as precise as the earlier ones, partly because of the incessant swaying of the towers.
The Michaelis-Thurm in Hamburg was never still all the time I was there. Its horizontal pendulum-like oscillations often exceed \(\frac{1}{2}\) arcminute.
Of all my triangulation points, this tower proved by far the worst as a station and (apart from the Brocken) also as a target.
\textsc{Repsold} also makes new bubble levels of extraordinary perfection, which he fills with naphtha instead of spirit of wine.
They come to rest much more quickly than the others.
